\section*{The Propers: Detailed Commentary}

\subsection*{The typical edition against the Benziger printing}

The Vatican \latin{editio typica} prints this formulary at pp.~397--398 under
marginal nos.\ 1592--1601. The Benziger \latin{editio iuxta typicam} of the same
year prints the same ten elements and differs from it in eight places.
\textbf{Two of the eight are substantive and six are pointing, abbreviation or
paging.}

\begin{receptiontable}{Place}{Vatican \latin{editio typica}}{Benziger \latin{iuxta typicam}}
Introit, after \latin{℣. Glória Patri.} & nothing follows & \latin{Miserére.} --- a repetition cue giving the antiphon's incipit. \textbf{Substantive.} The Benziger prints such cues habitually, \latin{Iustus.} at the Seventeenth Sunday on the same page image; the typical edition prints none here and none there\\
Secret, conclusion & \latin{Per Dóminum.} & \latin{Per Dóminum nostrum.} \textbf{Substantive}: the Benziger conclusion is the longer of the two\\
Introit verse & \latin{quóniam inops et pauper sum ego} & \latin{quóniam inops, et pauper sum ego} --- a comma. The 1861 Cummiskey's parallel Latin agrees with the typical edition; the 1862 Pustet agrees with the Benziger\\
Gospel, at \latin{humiliábitur} & \latin{humiliábitur: et qui se humíliat} & \latin{humiliábitur, et qui se humíliat}\\
Gospel, joining vv.~8--9 & \latin{ab illo, et véniens is qui te et illum vocávit} & \latin{ab illo, et véniens is, qui te, et illum vocávit}\\
Postcommunion, at \latin{páriter} & \latin{præsens páriter et futúrum} & \latin{præsens páriter, et futúrum}. The 1862 Pustet points it as the Benziger does and the 1604 sets a comma after \latin{præsens} instead: \textbf{the older books point this clause and the typical edition does not}\\
Heading style & \latin{DOMINICA / DECIMA SEXTA / post Pentecosten} on three lines; \latin{Ant. ad Introitum} and \latin{Ant. ad Offertorium} abbreviated, \latin{Antiphona ad Communionem} spelled out & two lines of capitals; \latin{Antiphona ad \dots} spelled out throughout\\
Running head & none for this Sunday anywhere in the book & \latin{Dominica XVI post Pentecosten} over its printed p.~392\\
\end{receptiontable}

\textbf{The typical edition is not uniform with itself inside this one
formulary}: it abbreviates two of the three antiphon headings and spells the
third out. And no page of the 1,088-page facsimile carries a running head
naming this Sunday --- printed p.~397 is headed \latin{Dominica XV post
Pentecosten} and p.~398 \latin{Dominica XVII post Pentecosten} --- so a finding
aid that locates a formulary by its running head reports this one absent.

Everything else agrees word for word and accent for accent. Both witnesses
print \latin{intellégimus} in the Epistle, both print the Communion citation as
\latin{Ps. 70, 16-17 et 18}, and both print \latin{éxtrahet} in the Gospel where
the controlling facsimile's own embedded text layer gives \latin{extranet}.

\textbf{The ten appointed texts do not move across the century before 1962.}
The Pustet Ratisbon of 1862, the Vatican typical edition of 1920, the Benziger
\latin{iuxta typicam} of the pre-1955 rubrics and the 1962 typical edition print
the same ten texts in the same order, and all four give the Alleluia as
\latin{Cantáte Dómino cánticum novum}. \textbf{The psalm citations of the two
middle books carry less than their texts do.} The optical layers of the 1920
typical edition and of the pre-1955 Benziger damage this formulary's own
numerals: the 1920 gives the Offertory's citation as \texttt{Ps. 39, 24 et 15},
where the 1962 book prints \latin{Ps. 39, 14 et 15}, and the Gospel's reference
as \texttt{huc. 24, 2-22}. \textbf{So what holds steady across the century
before 1962 is the ten texts and their order}, and the psalm citations of those
two books are not evidence of it. What changes around the texts is the
apparatus: the rank at the head is \latin{semiduplex} in 1920 and in
the Benziger and \latin{II classis} in 1962, and the three cross-references to
the seasonal supernumerary orations that stand in the two earlier books are gone
from the last.

\subsubsection*{Fifteen places where the appointed text is not the Bible's text}

The consequential collation is against the Clementine Vulgate, the Latin behind
the Douay--Rheims Challoner English printed with these propers. Fifteen
divergences stand between the two texts. \textbf{Four of the five chants depart
from the psalter and the Communion departs at no word at all.} The two books differ
besides in consonantal \latin{i} and \latin{j}, in the \latin{æ} and \latin{œ}
ligatures and in stress accents; those are the same words, spelled as each book
spells them.

\begin{divergencetable}
\cue{Int.} Ps. 85:3 & \latin{Miserére mihi, Dómine} & \latin{Miserere mei, Domine} --- a different word: the chant asks mercy \emph{for me} in the dative\\
\cue{Int.} Ps. 85:5 & \latin{quia tu, Dómine} & \latin{Quoniam tu, Domine} --- a different conjunction\\
\cue{Int.} Ps. 85:5 & \latin{suávis ac mitis es} & \latin{suavis et mitis} --- a different conjunction, and the chant supplies the verb the psalter leaves out\\
\cue{Int.} Ps. 85:5 & \latin{et copiósus in misericórdia} & \latin{et multæ misericordiæ} --- a different construction of the same sense\\
\cue{Int.} verse, Ps. 85:1 & \latin{aurem tuam mihi} & \latin{aurem tuam} --- the chant adds the dative; and the psalm title \latin{Oratio ipsi David} is not sung\\
\cue{Grad.} Ps. 101:16 & \latin{Timébunt gentes} & \latin{Et timebunt gentes} --- the chant drops the connective that opened the verse\\
\cue{Grad.} Ps. 101:17 & \latin{Quóniam ædificávit} & \latin{quia ædificavit} --- a different conjunction, the verb the same perfect in both\\
\cue{Grad.} Ps. 101:17 & \latin{vidébitur in maiestáte sua} & \latin{videbitur in gloria sua} --- a different noun, and the substantive divergence of this element\\
\cue{All.} Ps. 97:1 & \latin{quia mirabília fecit Dóminus} & \latin{quia mirabilia fecit} --- the chant names the subject the psalter leaves implicit; it stops before \latin{Salvavit sibi dextera ejus, et brachium sanctum ejus} and omits the title \latin{Psalmus ipsi David}\\
\cue{Off.} Ps. 39:14 & \latin{Dómine, in auxílium meum réspice} & \latin{Domine, ad adjuvandum me respice} --- a different construction of the same petition\\
\cue{Off.} Ps. 39:15 & \latin{confundántur et revereántur, qui quærunt} & \latin{Confundantur et revereantur simul, qui quærunt} --- the chant drops \latin{simul}; it also opens at the second half of v.~14, stops at the first half of v.~15, and repeats the second half of v.~14 as a refrain the psalm has not\\
\cue{Ep.} Eph. 3:13 & \latin{Fratres: Obsecro vos, ne deficiátis} & \latin{Propter quod peto ne deficiatis} --- the liturgical incipit adaptation\\
\cue{Ep.} Eph. 3:20 & \latin{quam pétimus, aut intellégimus} & \latin{quam petimus aut intelligimus} --- a different received spelling of one word, and not a different word\\
\cue{Gosp.} Luke 14:1 & \latin{In illo témpore: Cum intráret Iesus} & \latin{Et factum est cum intraret Jesus} --- the liturgical incipit adaptation\\
\cue{Gosp.} Luke 14:8--9 & \latin{ab illo, et véniens is qui te et illum vocávit} & \latin{ab illo. Et veniens is, qui te et illum vocavit} --- the missal runs the two verses into one sentence, lowercasing the \latin{Et} and setting a comma where the Clementine ends a sentence\\
\end{divergencetable}

Nothing else in the Gospel departs. \textbf{The Communion departs at nothing}:
every word is the Clementine's, and the antiphon's whole distance from the
psalter is where it begins and ends.

\textbf{At four of these five chants Cassiodorus's sixth-century lemma stands
with the missal against the Clementine.} He reads
\latin{suavis ac mitis es; et copiosus in misericordia omnibus invocantibus te}
at the Introit, \latin{videbitur in maiestate sua} at the Gradual, \latin{quia
mirabilia fecit Dominus} at the Alleluia and \latin{Domine, in auxilium meum
respice} at the Offertory; Migne's apparatus records no manuscript variant
against the Introit's clause or the Offertory's. At the Communion his text and
the Clementine both agree with the antiphon. \textbf{Migne's transcription is
not a critical edition}, so each of those four agreements rests on one word of
one printing. Theodoret's Greek names the Alleluia's subject too. \textbf{The Greek cuts the other way at two of the four}:
Theodoret's lemma at Ps.~101:17 is a future carrying the Clementine's noun and
not the chant's, and his Ps.~39:14 is the Septuagint's, which Migne's facing
Latin renders with the Clementine's \latin{ad adjuvandum me respice} and not
with the chant's clause.

\subsubsection*{Psalm numbering, and where the two modern series part}

\correction{Which psalm numbering the parentheses give}{The missal cites five
psalms in the Vulgate numbering, which the Douay carries: Ps.~85, 97, 101, 70
and 39 are Ps.~86, 98, 102, 71 and 40 in a modern Bible. \textbf{At two of the
five the modern verse number depends on which modern series is meant.} Ps.~39
and Ps.~101 carry their titles as numbered first verses, which English Bibles
leave unnumbered, so the Offertory is Hebrew 40:14--15 but English 40:13--14
and the Gradual is Hebrew 102:16--17 but English 102:15--16. Ps.~70, 85 and 97
fold their titles into verse one, so the Introit, the Alleluia and the Communion
answer to the same numbers in both series. \textbf{The same offset governs the
marginal figures on the New Advent pages of Augustine's \work{Enarrationes}},
which pair a Vulgate psalm number with a Hebrew verse number and therefore run
one behind the missal at Ps.~101 and Ps.~39. And the NPNF series titles each
exposition by the Hebrew psalm number, so Augustine on the Introit's Ps.~85 is
filed as ``Psalm 86''. \textbf{The printings of the three
commentators who stand beside him at these psalms depart from the missal's
numbering three further ways, and no two of the three are the same.} Migne's
numerals for Cassiodorus run level with the missal at Ps.~85, Ps.~101 and
Ps.~70, three ahead at Ps.~39, where they begin at
\latin{Exspectans exspectavi} and split two later verses, and one ahead from the
first verse at Ps.~97. The 1866 Bellarmine prints the appointed verses one
number \emph{below} the Vulgate at Ps.~39 and Ps.~101, whose superscriptions its printing omits, and
level with it at the other three. And in the Migne facsimile of Theodoret
pp.~878--879 reprint the same opening as pp.~876--877, so the two openings are
distinguished by column and not by page number.}

\subsection*{1. Introit \cue{Int.} --- \latin{Miserére mihi, Dómine}}

The antiphon is a centonisation of two verses that do not adjoin --- v.~3
entire, then v.~5 --- with the psalm verse drawn from v.~1. Nothing on the page
marks that v.~4 has been passed over. Its five departures from the psalter are
the most of any element here, and the commentators fall out two ways over the
antiphon's wording: Augustine's lemma differs from the chant at three places,
two of them among those five, and Cassiodorus's carries the chant's whole second
half.

\begin{witnessnote}
\wlocus{Augustine, \work{Enarratio in Ps.~85}, one sermon, read whole}
\S1 sets the whole psalm inside the \latin{totus Christus}: the Son of God
\latin{qui et oret pro nobis, et oret in nobis, et oretur a nobis}. \S2, on the
verse the Introit sings as its psalm verse: \latin{Inclinat aurem, si tu non
erigas cervicem: humiliato enim appropinquat; ab exaltato longe discedit, nisi
quem ipse humiliatum exaltaverit}, with the rich Pharisee boasting his merits
set against the needy publican. \S3 turns on who the poor man is, and answers
that the rich may be \latin{in Deo pauperes} (1 Tim.~6:17). \S5, on the
antiphon's own \latin{tota die}: \latin{omni tempore intellige \dots\ corpus
Christi tota die clamat, sibi decedentibus et succedentibus membris} --- the
cry is the Body's and runs through the members who die and the members who
follow. \S7, on \latin{suavis}: \latin{Quid est mitis? Portans me, donec
perficias me \dots\ nec meminit tantas quas incondite fundimus, et accipit unam
quam vix invenimus}. His lemma reads \latin{Miserere mei} where the chant sings
\latin{Miserére mihi} and \latin{egenus et inops ego sum} where it sings
\latin{inops et pauper sum ego}; at v.~5 it reads \latin{suavis es ac mitis},
nearer the chant than the Clementine. Latin from the Migne transcription and the
augustinus.it text; English from the \work{Nicene and Post-Nicene Fathers}, first
series, where the exposition is filed under the Hebrew number as ``Psalm 86''.
\end{witnessnote}

\begin{witnessnote}
\wlocus{Cassiodorus, \work{Expositio psalmorum}, in Ps.~LXXXV \S\S2, 4, 6, 8, 10, 18, 23--24}
He assigns all three appointed verses to Christ alone. His \latin{divisio} makes
the whole psalm Christ's prayer, \latin{in prima sectione dicens quae ipsi tantum
probantur aptari}, the section prayed \latin{pro membris suis} beginning only at
v.~11. At v.~1 the poor man is Christ in the form of a servant ---
\latin{Inclina, ostendit quia se ad ipsum extendere non poterat humana conditio,
nisi ipse suam Maiestatem piissimus inclinaret} --- and the poverty is
\latin{conditionem, quam susceperat humanitatis, quae ex se nihil habere potest,
nisi quod largitate Divinitatis acceperit}; he turns it pastorally at once,
\latin{Audiant egeni et pauperes}, and does not shut out the rich,
\latin{egenus enim et pauper Dei est, quisquis meruerit mundi istius perversitate
vacuari}. \textbf{His \latin{tota die} is one human life and not the Body's
succession}: \latin{totius vitae tempus ostenditur, ut per multa tempora
annorumque curricula quasi unius diei continuus clamor esse monstretur}, and the
crying is done by works, \latin{Magna siquidem voce ille clamat ad Dominum, qui
quamvis lingua taceat, bonis tamen operibus perseveranter exclamat}. At v.~5 he
takes \latin{suavis ac mitis} as a formal definition, \latin{per quintam speciem
definitionis, quae Graece dicitur kata ten lexin}: \latin{Suavis, quia post
amaritudinem huius mundi dulcis est ad se recurrentibus. Mitis, quia diu sustinet
peccatores. Copiosus in misericordia, quia licet sint nostra numerosa peccata,
multo abundantior est pietas.} \textbf{His lemma there is the antiphon's second
half word for word}, and his v.~1 reads \latin{aurem tuam ad me} against the
missal's \latin{mihi} and \latin{egenus} against \latin{inops}.
\end{witnessnote}

\begin{witnessnote}
\wlocus{Theodoret of Cyrus, \work{Interpretatio in Psalmos}, PG~80 cols.\ 1553--1556}
Antiochene throughout. His hypothesis is historical --- \latin{Praecinit autem et
Assyriorum in Hierosolymam impetum, et Ezechiae spem} --- and the poverty is a
philosopher's and not a moralist's: \latin{Nam uterque justitiae divitias
possidens, et divinus David, et admirabilis Ezechias, has quidem non
considerabant, ad naturae vero paupertatem respiciebant}. \latin{Inclina, Domine,
aurem tuam} he takes \latin{per metaphoram desumptam ab aegroto, qui ob
imbecillitatem non potest clare loqui, et medicum cogit, ut aurem ori ejus
admoveat}. \textbf{At v.~3 he allegorises \latin{tota die} not at all} and
reports a version instead: \latin{Nam toto die, Symmachus singulis diebus dixit.}
At v.~5 he supplies the Greek the Latins argue from without, and reports what
the other versions made of the antiphon's second adjective: \latin{Pro mitis
autem Aquila et Theodotio placabilis dixerunt. Mansuetudo igitur patientiam
quoque innuit.} His Migne Latin at v.~1 reads \latin{inops et pauper}, which is
the missal's pair where Cassiodorus's is not. This is Migne's reprint of
Schulze's recension with a facing Latin version, not a critical Greek text, and
the Greek and that Latin do not always agree.
\end{witnessnote}

\begin{witnessnote}
\wlocus{Robert Bellarmine, \work{Explanatio in Psalmos}, in the 1866 English of John O'Sullivan, printed pp.~259--261}
The form of the prayer is his point at v.~1: ``He begins his prayer by touching
on God's greatness and his own poverty, an excellent form of prayer \dots\ As I
am the beggar sitting at the rich man's gate, incline thy ear to your poor
servant.'' The poor man is a disposition and not a purse --- ``the person, who,
though he may abound in the riches of the world, still does not put his trust in
them, takes no pride in them'' --- and he cites Augustine that ``Lazarus was not
taken up into Abraham's bosom by reason of his poverty, but on account of his
humility.'' At v.~3 \latin{tota die} is perseverance in prayer. At v.~5 he
glosses each word of the antiphon in turn, attributes the distracted-prayer
reading of \latin{mitis} to Augustine, and extends it with his own image of the
judge and the culprit who turns aside to talk with his friends. O'Sullivan's
translation is an abridgement its own translator declares, and the argument line
and psalm text printed above each exposition in this edition are the Douay's and
not Bellarmine's.
\end{witnessnote}

\textbf{All four read \latin{inops et pauper} as the confession of a man who has
nothing of his own, and none of them reads it as a statement about money}; two
of them volunteer the qualification that the rich are not excluded, Cassiodorus
and Bellarmine, and Augustine reaches it at \S3's question of who the poor man
is, answering that the rich may be \latin{in Deo pauperes}. Each supplies a different image for the bending of the
ear --- Theodoret's sick man too weak to speak, whom
the physician must stoop to; Cassiodorus's human condition that cannot stretch
upward; Bellarmine's beggar at the rich man's gate; Augustine's ear bent only to
the neck that is not stiffened. \textbf{And they divide four ways on \latin{tota
die}}: Augustine's is the whole time of the Body, Cassiodorus's a single human
life crying by its works, Bellarmine's fervour and perseverance in prayer, and
Theodoret's Symmachus's ``on each day'' and nothing further. An argument that
leans on the singularity of one long day has the fourth of them against it.

\textbf{Aquinas does not reach this psalm.} His \work{Postilla super
Psalmos} as the Corpus Thomisticum serves it breaks off inside Psalm~54, so of
this Mass's five psalms he is a direct commentator on the Offertory's alone.

Four canonical relations belong to the appointed words. \latin{Mitis} as a whole
word stands in exactly two verses of the Clementine, Ps.~85:5 and Matt.~11:29,
where Christ pairs it with \latin{humilis corde}; \latin{iugum meum suave est}
follows immediately at 11:30, \textbf{so it is across those two verses that the
antiphon's \latin{suávis} and \latin{mitis} are joined}, and neither verse
carries both alone. \latin{Copiósus in misericórdia} answers no phrase anywhere
in that Bible; the psalter's own formula stands at Ps.~85:5 itself, at
Ps.~85:15, at Num.~14:18 and at Joel~2:13. The dative the chant adds gives
Ps.~85:1 the petition-form \latin{inclina aurem tuam mihi}, which in that
psalter stands at Ps.~16:6 and nowhere else, against the \latin{inclina ad me
aurem tuam} of Ps.~30:3, Ps.~70:2 and Ps.~101:3 --- the last two the psalms of
this Mass's own Communion and Gradual, and not one of the three carrying
\latin{mihi} at all. \textbf{What the two verses share is that four-word
petition and the \latin{et exáudi} beside it, and nothing wider}: the Introit
sings \latin{Inclína, Dómine, aurem tuam mihi, et exáudi me}, where Ps.~16:6
carries no vocative, asks that God hear \latin{verba mea}, and opens on a clause
the psalm verse has not. And Apoc.~15:4
brings together the nations coming to adore and the fear of the divine name,
which are Ps.~85:9 and Ps.~101:16, though Ps.~85:9 is not appointed here.

\textbf{Four psalms in the whole Clementine open with \latin{Oratio}, and this
Mass appoints two of them} --- Ps.~85, \latin{Oratio ipsi David}, and Ps.~101,
\latin{Oratio pauperis} --- while a third, Ps.~16, \latin{Oratio David}, is the
psalm whose v.~6 the added dative reproduces. Neither title is sung: the missal
omits Ps.~85's, and Ps.~101's stands fifteen verses before the Gradual's excerpt.

No New Testament book quotes Ps.~85:1, 3 or 5, and the antiphon's English
produced no later reuse: \latin{Miserere mihi} is not the famous
\work{Miserere}, which is Ps.~50, and the secular lives of that word in English
attach there.

\subsection*{2. Collect \cue{Coll.} --- \latin{Tua nos, quǽsumus, Dómine, grátia}}

One grace is asked to do two things and then to produce a third: to go before,
to follow, and to make us continually intent on good works. \textbf{No Father
expounds this prayer}, which is the ordinary state of a Roman sacramentary
oration: no commentator working through a Sunday's own orations stands behind
it. \textbf{The one witness who
comments on this Collect as this Mass's own is a nineteenth-century liturgist},
and he writes for devotion rather than argument.

\begin{witnessnote}
\wlocus{Prosper Guéranger, \work{The Liturgical Year}, vol.~XI, printed p.~357}
Above the Latin of this Collect and its English he sets the distinction the two
verbs make: ``unless grace prevent, that is, anticipate, us, we cannot have so
much as the thought of doing what is holy; and again, unless it follow up the
inspirations it has given us, and lead them to a happy termination, we shall
never be able to pass from the simple thought to the act of any virtue
whatsoever. If, on the other hand, we be faithful to grace, our life will be one
uninterrupted tissue of good works.'' He sets \emph{prevent}, \emph{follow},
\emph{grace} and \emph{good works} in italic on this page; the word
\emph{prevent} stands in this exposition and not in his own English of the
prayer, which renders \latin{prævéniat} as ``ever go before us.''
\textbf{He witnesses that the Collect is read this way and is not the
authority for the technical distinction}, which is Augustine's and Trent's; and
his ``persevering continuity of this most precious aid'' is his phrase and not a
rendering of \latin{iúgiter}.
\end{witnessnote}

\begin{witnessnote}
\wlocus{Augustine, \work{Enchiridion} 32 and \work{De gratia et libero
arbitrio} 17.33} \latin{Misericordia eius praeveniet me \dots\ Misericordia eius
subsequetur me: nolentem praevenit, ut velit, volentem subsequitur, ne frustra
velit.} The scriptural substratum is Ps.~58:11, which stands verbatim in the
Clementine; the following clause answers to Ps.~22:6, where the Clementine reads
\latin{misericordia tua subsequetur me} against Augustine's \latin{eius}. In
\work{De gratia et libero arbitrio} the same pair is put as \latin{operatur} and
\latin{cooperatur}.
\end{witnessnote}

\begin{witnessnote}
\wlocus{Council of Trent, Session~VI} Cap.~XVI: \latin{virtus \dots\ bona eorum
opera semper antecedit et comitatur et subsequitur}, which shares
\latin{iugiter} and \latin{semper} with this Collect at the doctrinal point.
Cap.~X and can.~32 stand with it. \textbf{Trent's three verbs are not the
Collect's two}: \latin{comitatur} is the council's addition. Cap.~X does quote a
Sunday collect of this series for its doctrine, and the Tauchnitz printing's own
footnote identifies it as the Thirteenth Sunday's --- quoted altered. So a
council does cite a collect of this series, and it does not cite this one.
\end{witnessnote}

\textbf{A second Latin voice reaches the prevenient half of the doctrine at an
appointed text of this Mass.} Cassiodorus opens his exposition of the
Communion's psalm on Christ's charity \latin{quae nullis meritis praecedentibus
gratis semper impenditur} and closes it \latin{Cum totus hic psalmus gratiam
Domini, quae gratis datur, summa intentione commendet}; at the Offertory's psalm
he makes God's very looking our protection, \latin{aliter enim liberari non
possumus, nisi nos Divinitas propitiata respiciat}. \textbf{The first of those
says at that same psalm what Augustine says there in \latin{nihil tuum
praecessit} and \latin{docuisti me nihil in me praecessisse}}, \work{Enarratio
in Ps.~70} sermo~II \S\S1--2, and Cassiodorus names \latin{doctor Augustinus}
inside these same five psalms, so at that clause the two Latins are one judgment
restated and not two.

The Collect's \latin{bonis opéribus} and the Postcommunion's \latin{rénova
\dots\ mentes nostras} share vocabulary with Eph.~2:10 and Eph.~4:23, both
outside the appointed pericope; those are echoes, and the three orations of this
Mass quote no Scripture at all.

\subsection*{3. Epistle \cue{Ep.} --- Eph. 3:13--21}

\textbf{The lesson begins one verse before the sentence it completes.} Eph.~3:1
opens \latin{Huius rei gratia} and breaks off into a digression that runs to
v.~13; v.~14 resumes with the same three words. What a congregation hears is the
resumption without the thing resumed, and the liturgical incipit removes even
that signal by substituting \latin{Fratres: Obsecro vos} for the Clementine's
\latin{Propter quod peto}. At the other end the pericope closes the doctrinal
half of the letter: the exhortation opens at 4:1 with \latin{Obsecro itaque
vos}, which is the same verb the liturgical incipit put at the front --- a
coincidence of the two Latin texts and not a derivation. Every petition inside
the reading asks that God give, and every principal noun of the petition is a
word Ephesians has already used.

\textbf{The calendar registry named in the References gives Ephesians~3 three
appointments and two of them cut v.~13 out}, the Sacred Heart and St~Margaret
Mary Alacoque reading 3:8--12, 14--19 and 3:8--9, 14--19; and in that registry
\textbf{vv.~20--21, the doxology, stand here and at no other Mass}.

\begin{witnessnote}
\wlocus{John Chrysostom, \work{Homily 7 on Ephesians}} The homily runs from 3:8
to the doxology of 3:21, so the whole appointed pericope stands inside it and
Homilies 6 and 8 fall outside it on either side. At v.~13 he answers the
difficulty in the Apostle's own terms: the tribulations are the readers'
glory because God so loved them as to afflict his servants for them, with
Hos.~6:5. At vv.~14--15 the bowing of the knees shows that the prayer is
heartfelt, and \latin{omnis patérnitas} means that the tribes are now reckoned
by their Creator rather than, as before, ``according to the number of Angels.''
\textbf{At the four measures he does not read the Cross}: they are the immensity
of God's love ``and how it extends every where,'' sketched by the visible
dimensions of a solid body, ``pointing as it were to a man.'' Chrysostom is held
here in Gross Alexander's English in the \work{Nicene and Post-Nicene Fathers},
first series, vol.~13, and in no Greek text.
\end{witnessnote}

\begin{witnessnote}
\wlocus{Augustine, \work{Sermo}~165} Headed \latin{De verbis Apostoli (Eph 3,
13-18) \dots\ deque gratia et libera voluntate, contra Pelagianos}. \S1.1 reads
the Epistle's hinge as a distinction between what the hearer's will can do and
what it cannot: \latin{Quia ergo voluntatis habetis arbitrium: Peto. Quia vero
voluntatis non sufficit arbitrium ad implendum quod peto: Huius rei gratia
flecto genua mea ad Patrem \dots\ Peto enim a vobis, propter arbitrium
voluntatis: rogo det vobis, propter auxilium maiestatis.} \S\S3.3--5.5 assign
the four measures; \S7.9 closes on \latin{unde gratiam tuam meritum meum
praecedat? Non}. \textbf{The sermon covers 3:13--18 and
stops three verses short of the doxology}, and its opening remark that the
day's Apostle, psalm and Gospel all agreed \latin{ut spem non in nobis, sed in
Domino collocemus} is evidence about Augustine's own African order --- the psalm
sung there was Ps.~56:2 --- and says nothing about this Roman formulary. At
\S2.2 he pauses to defend the translator's \latin{altitudo}, which is his
African text's reading where the missal and the Clementine print
\latin{sublímitas}: \textbf{a sentence of his about \latin{altitudo} is not a
sentence about the missal's word.}
\end{witnessnote}

\begin{witnessnote}
\wlocus{Augustine on the four dimensions elsewhere: \work{Epistula}~55.14.25 and
\work{In Iohannis euangelium tractatus}~118.5} Both name Eph.~3:18 as the text
expounded, so neither attribution is an inferred echo, and both read the four
measures as the four members of the Cross. \textbf{They differ from each other
on two of the four assignments and reverse the order of height and length}, and
\work{Sermo}~165 differs again --- breadth the charity that alone works well,
length perseverance to the end, height the \latin{sursum cor} and the loving of
God \latin{gratis}, depth the unsearchable judgments of God. \textbf{Three
cross-readings with three sets of virtues is a development inside one author},
and no one of them is ``the'' patristic reading of the verse. English from the
\work{Nicene and Post-Nicene Fathers}, and the Latin from Migne rather than
from CCSL~36 or CSEL~34.
\end{witnessnote}

\begin{witnessnote}
\wlocus{Thomas Aquinas, \work{Super Epistolam ad Ephesios lectura} cap.~3
lect.~4--5} Lect.~4 treats v.~13. At lect.~5 he sets out three readings of the
measures and chooses none: the Dionysian, taking the dimensions of God
metaphorically and grounded on Job~11:8--9; the charity reading; and the Cross,
where he specifies the hidden depth as \textbf{predestination}. \textbf{He
names Augustine nowhere in the cross-reading.} Marietti text, not Leonine.
\end{witnessnote}

\begin{witnessnote}
\wlocus{Prosper Guéranger, \work{The Liturgical Year}, vol.~XI, printed pp.~358--360}
He takes the four measures neither of the Cross nor of the extent of love, but of
the indwelling: ``God alone \dots\ can strengthen in us the inward man enough to
make us understand, as the saints do, the dimensions (breadth, length, height,
and depth) of the great mystery of Christ dwelling in man, and dwelling in him
for the purpose of filling him with the plenitude of God.'' And he reads the
lesson as this Mass's Collect prays: ``Now, it depends on us to follow God's
grace; nothing else but our own resistance prevents the Holy Ghost from making
saints of us.'' He observes at p.~358 that the Church borrows more from Ephesians
in this season than from any other Pauline letter.
\end{witnessnote}

\textbf{So the four measures are read at least three ways by the witnesses
standing over this verse} --- the extent of love, the members of the Cross, the
dimensions of the indwelling --- and the cruciform reading is one of them and not
the reading. The Antwerp 1614
\work{Commentaria in omnes D. Pauli Epistolas} of Cornelius a Lapide, whose title
promises all the epistles, \textbf{stops at Galatians 6:17}: no \latin{ad
Ephesios} running head stands anywhere in it. Lapide did comment on Ephesians;
that book does not carry it.

Two Pauline places carry the reading's own vocabulary from outside the pericope.
2 Cor.~4:16--17 puts \latin{non deficimus}, the inward renewal, tribulation and
a weight of glory in one movement; Col.~1:24 is the nearest Pauline statement of
the exchange Eph.~3:13 makes between an apostle's affliction and his hearers'
glory.

\subsection*{4. Gradual \cue{Grad.} --- \latin{Timébunt gentes nomen tuum}}

The two verses stand at the psalm's hinge, where it turns from the wasting of
one afflicted man to the rebuilding of Sion; the psalm's own title makes the
whole a poor man's prayer. Four commentators expound them, and they divide on
the tense, on the noun, on who is speaking and on what is seen.

\begin{witnessnote}
\wlocus{Augustine, \work{Enarratio in Ps.~101}, sermo~I \S\S16--17} \S16 reads
the fear of the name as the ingathering of the Gentiles, who become the second
wall meeting Israel in the corner stone --- and he cites \textbf{Eph.~2:20},
a verse of the appointed Epistle's own letter which this Mass does not read.
\S17 reads a future: Sion is being built now, \latin{Hoc agitur nunc}, and the
seeing is the Judgment, set against the first coming when he was seen \latin{in
infirmitate sua} and, in the \work{Nicene and Post-Nicene Fathers} English,
``had no form nor comeliness'' (Isa.~53:2). \S18 expounds
the verse the Gradual stops before, \latin{Respexit in orationem humilium}, and
\S19, three verses past the Gradual's cut, gives \latin{ex alto factus est
humilis, ut humiles exaltaret}. Both sections fall inside Sermo~I, before the
heading \latin{SERMO II. De secunda parte Psalmi}.
\end{witnessnote}

\begin{witnessnote}
\wlocus{Cassiodorus, \work{Expositio psalmorum}, in Ps.~CI \S\S2--3, 22--23}
\textbf{His lemma is the missal's at both disputed words}: \latin{Quia
aedificavit Dominus Sion, et videbitur in maiestate sua}. At v.~16 the fear of
the nations is historical and ascetic rather than ecclesiological: the true Lord
was not feared while the world served idols, and after the saving advent
\latin{gentes conversae sunt per timorem}; the kings of the earth are those who
bridled their own bodies, \latin{qui corpora sua divinis regulis infrenantes, sui
imperatores esse (Domino praestante) valuerunt}. At v.~17 the building is
finished and the thing built is the Church, \latin{hoc est mater Ecclesia, de
vivis lapidibus fabricata, in qua Domini cultura usque ad finem mundi sine
intermissione proficiet}, and the seeing in majesty is the Judgment in the
assumed body, \latin{quando haedos sequestrat ab agnis}. \textbf{And he records
Augustine's reading of the speaker and declines it} (\S2): \latin{Quamvis aliqui
praesentem psalmum Domini Salvatori aptandum esse putaverint, conveniens tamen
videtur afflicti magis et gementis pauperis \dots\ quia multa sunt quae illi
immaculatae sanctae incarnationi nequeunt convenire. Et primum, quod anxius
nusquam fuisse legitur Dominus Christus} --- assigning the psalm instead to an
afflicted penitent poor man left nameless on purpose, \latin{ut cum uni datur,
omnes sibi pauperes Christi cognoscerent attributum}.
\end{witnessnote}

\begin{witnessnote}
\wlocus{Theodoret of Cyrus, \work{Interpretatio in Psalmos}, PG~80 cols.\ 1679--1682}
\textbf{His Greek lemma at v.~17 is a future, and its noun is the Clementine's
and not the chant's}, where Migne's facing Latin quietly prints the Vulgate's
perfect \latin{aedificavit} on the opposite column. \textbf{And his exposition is historical, not
eschatological}: the rebuilding of the city is itself God's glory and answers the
charge that the exile proved him weak, \latin{Sione rursus aedificata, in
pristina gloria universorum Deum omnes cernent}. At v.~16 he holds the historical
and the christological together and calls the fulfilment partial, \latin{hoc vero
proprie ac vere post Dei et Salvatoris nostri incarnationem contigit}, with an
explicit already-and-not-yet on Heb.~2:8 and Phil.~2:10 --- \latin{omne genu ipsi
flectendum esse}. \textbf{The speaker, for him, is the people in the Babylonian
captivity.}
\end{witnessnote}

\begin{witnessnote}
\wlocus{Robert Bellarmine, \work{Explanatio in Psalmos}, printed p.~317, his own vv.~15--16}
He holds together what Augustine and the missal divide. The building is
accomplished ``in the present day, having established his Church in spite of all
kings and nations,'' and the seeing is ``in the time to come, when he shall come
with all his Angels \dots\ When he began to build up Sion he was seen in his
lowliness. `We have seen him, and there was no sightliness' \dots\ but when he
shall come to pass judgment, then `he shall be seen in his glory.''' \textbf{That
antithesis is Augustine's own, made from the same Isaian verse.} At the preceding
verse he glosses the psalm's ``thy glory'' as ``that is, thy \emph{majesty}'' ---
a gloss on the sense, his own printed lemma reading ``in his glory.''
\end{witnessnote}

\textbf{The chant's two disputed words divide the witnesses along two different
seams.} The tense divides by text tradition: the Greek Theodoret expounds and
Augustine's Latin both read a future, while the Clementine, Cassiodorus and
Bellarmine read the perfect the chant prints. The noun divides Latin from Latin:
the chant and Cassiodorus read \latin{maiestáte} where the Clementine, Augustine
and Bellarmine's English read \latin{gloria}. \textbf{So no sentence built on
\latin{maiestáte} is Augustine's}, and the difference between him and the chant
is not a peculiarity of the missal.

\textbf{The exegetical divergence is sharper than either, and it does not follow
the tense.} All three Latins refer \latin{vidébitur} to the Judgment; Theodoret
refers it to Jerusalem restored. And Guéranger, above this Gradual at printed
p.~363, refers it to neither: ``The Church, which is showing herself in the midst
of the Gentiles, bears on herself the mark of her divine Architect; God shows
Himself, in her, \emph{in all majesty}; and, by her, the kings of the earth are
made to fear Him'' --- the majesty seen now, in the Church. He is introducing a
chant rather than construing a psalm, and the eschatological reading is not the
whole Latin position once he is counted.

\textbf{The verse the Gradual stops before divides the four again, three ways.}
Augustine and Cassiodorus take the poor whose prayer is regarded for the
faithful of Christ, one and many; Bellarmine refers them to the martyrs of
Apoc.~6:10; Theodoret to the captives whose prayer was admitted, \latin{Non enim
despexit, veluti captivos et servos}. What follows the chant's cut is the
psalm's and not the chant's.

Heb.~1:10--12 quotes this psalm of the Son, but at vv.~26--28, outside the
appointed excerpt; the Epistle's \latin{radicáti, et fundáti} shares its verb
with \latin{fundasti} at that same v.~26, nine verses past the Gradual, and the
Gradual's own building word is \latin{ædificávit}. These two verses produced no
later reuse: ``the Gentiles shall fear thy name'' returns no page at all in the
Library of Congress's \work{Chronicling America}.

The Gradual is not proper to this Mass. \textbf{The calendar registry named in the
References} also gives it to the Third, Fourth, Fifth and Sixth Sundays after the Epiphany, and gives its
first verse as the Alleluia verse of the Eighteenth Sunday, two Sundays after
this one; Schuster names the Epiphany Gradual independently. \textbf{It does not,
however, return with the resumed Epiphany Sundays in November}: a resumed Sunday
keeps its own orations, Epistle and Gospel and borrows the Twenty-third Sunday's
chants, a rule \work{The 1962 Roman Calendar: Rules, Ranks, and Reckoning}
prints from \latin{Rubricae generales} n.~18 and Missal rubric 298.

\subsection*{5. Alleluia \cue{All.} --- \latin{Cantáte Dómino cánticum novum}}

The chant prints two clauses and stops, and the four commentators who stand over
them give four answers to the one question the clauses leave open: what the
\latin{mirabília} are.

\begin{witnessnote}
\wlocus{Augustine, \work{Enarratio in Ps.~97}, \S1} Nothing later in the psalm
is appointed, so the exposition of the appointed words is \S1 alone. The new
song belongs to the new man; and because the psalm bids the whole round of the
earth sing it, those who cut themselves off from the communion of the whole
earth cannot sing it --- \latin{quia canticum novum in toto, non in parte
cantatur}. The \latin{mirabilia} are the raising of the whole world from
everlasting death by the Lord's holy arm, which he identifies with Christ out of
Isa.~53:1. And the passage moves from the new song to the new man to an inward
healing: \latin{Quis est qui intus sanatur? Qui credit in eum \dots\ in novum
hominem reformatus.} His lemma for the second half of the
verse reads \latin{Sanavit ei dextera eius} --- ``healed,'' and ``for him'' ---
and he builds a distinction on it, \latin{Multi enim sanantur sibi, non ei},
where the Clementine reads \latin{Salvavit sibi} and Cassiodorus a third form,
\latin{salvavit eum}; the Alleluia stops before that clause entirely. The NPNF
English carries its editors' own ellipses inside the sections, so the Latin is
the fuller witness; even there his treatment of this verse is short.
\end{witnessnote}

\begin{witnessnote}
\wlocus{Cassiodorus, \work{Expositio psalmorum}, in Ps.~XCVII \S\S2, 4, 6--7, 14}
\textbf{His lemma carries the chant's \latin{Dominus} against the Clementine.}
The new song is regeneration and the Incarnation, and he gives it the
sacramental term this Mass's Postcommunion uses: \latin{Propheta fideles admonet
Christianos, ut novae regenerationis sacramenta sumentes, novum canticum de
Domini incarnatione concelebrent. Novus enim homo cantare debet canticum novum,
non ille vetustus qui necdum Adae peccata deponens, in praevaricatione veteris
hominis perseverat.} On the wonders he first takes the Gospel healings ---
\latin{quando caecis lumen, claudis gressum, surdis etiam donavit auditum} ---
and then denies at once that these are the singular thing, \latin{Sed ista
fecerunt et sancti eius}, locating the novelty in the Resurrection.
\end{witnessnote}

\begin{witnessnote}
\wlocus{Theodoret of Cyrus, \work{Interpretatio in Psalmos}, PG~80 cols.\ 1657--1658}
The most austere of the four: the wonders are simply what exceeds nature and
expectation, \latin{Supra naturam enim et conceptum sunt ea quae a Deo
universorum fiunt}, and the new song answers a change of
dispensation --- his Greek word is one for a polity or manner of common life,
which Migne's facing Latin flattens to \latin{religionem}. \textbf{His lemma too
names the Lord}, the subject the chant supplies and the Clementine leaves
implicit. His rule that in this idiom \latin{Manum actionem appellari,
Dexteram vero, bonam actionem} is a lexical principle and not an allegory.
\end{witnessnote}

\begin{witnessnote}
\wlocus{Robert Bellarmine, \work{Explanatio in Psalmos}, printed p.~306}
He answers the same question by enumeration, and the list gathers both of the
readings the Fathers give apart: conceived of the Holy Ghost, born of a virgin,
committed no sin, justified sinners, the deaf hearing and the blind seeing, the
dead raised, ``and, what is the most strange and wonderful of all, shewed himself
alive within three days after he was buried \dots\ and, finally, as St.~Augustine
says, conquered the world, not by the sword but by the cross.'' He and Theodoret
agree, in almost the same words, that the psalm foretells both advents and dwells
on the first.
\end{witnessnote}

The chant supplies \latin{Dóminus} as the subject of \latin{fecit} and stops
before the Gentile horizon of vv.~2--3. \latin{Canticum novum} occurs in the
Clementine at Ps.~32:3, 39:4, 95:1, 97:1, 143:9, 149:1, Isa.~42:10,
Apoc.~5:9 and 14:3, and that list is exhaustive; four of the nine carry the
chant's full incipit \latin{Cantate Domino canticum novum} --- Ps.~95:1,
Ps.~97:1, Ps.~149:1 and Isa.~42:10. \textbf{Two of those four touch
this formulary's other elements}: Ps.~95:1 stands under the title \latin{quando
domus aedificabatur post captivitatem}, beside the Gradual's \latin{ædificávit
Dóminus Sion}; and Ps.~149:1 continues \latin{laus eius in ecclesia sanctorum},
beside the Epistle's \latin{cum ómnibus sanctis} and \latin{glória in
Ecclésia}. A third place in the list of nine, Ps.~39:4, belongs to the
Offertory's own psalm.

\textbf{The Gradual's psalm itself reaches for this verse two verses past the
Gradual's cut.} At Ps.~101:19 Cassiodorus quotes the Alleluia outright ---
\latin{Hic laudabit Dominum, novo scilicet cantico, sicut in alio psalmo dictum
est: Cantate Domino canticum novum} --- while Augustine and Theodoret reach
there for 2~Cor.~5:17 and the new creation. The join is inside that psalm; this
Mass sings the new song immediately after the Gradual.

The chant's full incipit is not distinctive to this verse, standing also at
Ps.~95:1, Ps.~149:1 and Isa.~42:10, and no later use of it can be tied to
Ps.~97:1.

\subsection*{6. Gospel \cue{Gosp.} --- Luke 14:1--11}

\textbf{The pericope cuts one continuous sabbath meal in half.} It keeps the
healing and the parable of places and stops before the instruction to invite the
poor and before the parable of the great supper; the bread motif frames the
whole scene at v.~1 and v.~15, and only its first half is appointed. The sabbath
argument repeats, with the same pair of animals, the argument of Luke~13:15,
where Matthew's analogue at 12:11 argues from one sheep, and this is the third
and last of Luke's sabbath-healing controversies. \textbf{The calendar registry
gives Luke~14 in three separated blocks} --- 14:1--11 here, 14:16--24 at the
Second Sunday after Pentecost, and 14:26--33 and 14:26--35 in the sanctoral and
a Common --- and it appoints vv.~12--15 and v.~25 nowhere.

Four commentators expound the pericope directly and they divide over what its
first half is about.

\begin{witnessnote}
\wlocus{Cyril of Alexandria, \work{Commentary on Luke}, Sermons 101 and 102}
Two consecutive sermons dividing the pericope exactly where it divides
internally. Sermon 101 makes the healing an argument about the sabbath: the law
was shadow and type waiting for the truth, keeping the sabbath rationally means
ceasing from sins and not from mercy, and the lawyers' silence is malice ---
pressed to sarcasm, that a man had better commit his child joyfully to the grave
than break the day. Sermon 102 reads the seating parable as the ordinary virtue
of a modest mind: to seize an honour not due is like a theft, with the stolen
goods to be given back, and worldly honour is grass on a housetop.
\textbf{Cyril's commentary survives entire only in Syriac}, and the English here
is R.~Payne Smith's 1859 rendering of that Syriac, so the words a reader meets
are a translator's twice over. \textbf{His lemma at 14:5 reads ``whose son of
yours, or whose ox''}, and the rhetorical force he draws from a father's
endangered child does not transfer to the missal's \latin{ásinus}.
\end{witnessnote}

\begin{witnessnote}
\wlocus{Ambrose, \work{Expositio in Lucam} VII.195--196} The whole pericope gets
one paragraph. The dropsical man is one \latin{in quo fluxus carnis exuberans,
animae gravabat officia, spiritus exstinguebat ardorem}; humility follows as the
next lesson, taught by forbidding the appetite for the higher place and taught
gently, \latin{ut humanitas persuasionis coercitionis asperitatem excluderet};
then hospitality, which is the Lord's kind only when spent on the poor,
\latin{nam hospitalem remuneraturis esse affectus avaritiae est}. \textbf{He
passes over the sabbath controversy that occupies Cyril's whole first sermon},
and books VIII and IX carry no second treatment of the chapter: his own heading
in book VII reads \latin{Cap. XIV. --- Vers. 2--14}, and the paragraph is the
whole of what he gives it.
\end{witnessnote}

\begin{witnessnote}
\wlocus{Augustine, \work{Quaestiones euangeliorum} II q.~29} The dropsical man
is compared to the animal fallen into the well, because he laboured under a
humour, and then to the covetous rich man. The numbering here is Migne's and not
CCSL~44B's.
\end{witnessnote}

\begin{witnessnote}
\wlocus{Bede, \work{In Lucae euangelium expositio}, PL~92 cols.\ 510D--513A}
The fullest continuous Latin exposition, and the load-bearing one for this
formulary. He argues the sequence from the body's disease to the heart's ---
\latin{per huius aegritudinem corporis, in illis exprimeretur aegritudo cordis}
(511B) --- after setting out the etymology and pathology of dropsy (511A). The
ox and the ass are the wise and the dull, or the two peoples, all found sunk in
the pit of concupiscence and drawn out \latin{iustificati gratis per gratiam
ipsius}, Rom.~3:24 (511C). The marriage feast is the union of Christ and the
Church, and the first place is glorying in one's own merits: \latin{non se, de
meritis gloriando, quasi caeteris sublimior extollat} (512A) --- which is the
vocabulary of this Mass's Collect and Communion, and is why he is the witness
that most nearly meets the formulary. \textbf{At \latin{Tunc erit tibi gloria}
he holds two readings and does not choose}: eschatological, \latin{ne nunc
quaerere incipias quod tibi servatur in fine}, and present, \latin{Potest autem
et in hac vita intelligi \dots\ quia quotidie Dominus suas nuptias intrat}
(512D). At v.~11 he argues that the conclusion proves the parable must be read
\latin{typice}, because it is plainly false that everyone who exalts himself
before men is humbled by men: \latin{omnis qui se incaute de meritis allevat,
humiliabitur a Domino} (513A). In the answering letter to Acca prefixed to this
work he says he compiled it from Ambrose, Augustine, Gregory and Jerome and
marked each borrowing with the author's initial in the margin; \textbf{the
Migne transmission carries no such marks}, and the dropsy-and-miser comparison at 511A stands almost word for word
in Augustine's \work{Quaestiones}, which the \work{Catena aurea} attributes to
Augustine.
\end{witnessnote}

\textbf{The \work{Catena aurea in Lucam} cap.~14 lect.~1--2 follows the pericope
exactly and names eight authorities.} Four of them --- Cyril, Bede, Ambrose and
Augustine's \work{Quaestiones} --- stand above at their own works.
Theophylact, Gregory's \work{Moralia}, Basil on the order of places at table,
and a Chrysostom excerpt at lect.~2 stand there only in the Catena's medieval
Latin of the Greek Fathers, which is not their Greek.

The closing sentence has a reception the pericope does not. \textbf{Benedict
opens \work{Regula}~7, \latin{De humilitate}, by quoting Luke~14:11 as the voice
of Scripture itself} --- \latin{Clamat nobis Scriptura divina, fratres, dicens:
Omnis qui se exaltat humiliabitur et qui se humiliat exaltabitur} --- draws from
it that every self-exaltation is a species of pride, and builds on it Jacob's
ladder and the twelve degrees. \textbf{He names no chapter, and the sentence
stands three times in the canon}, at Luke~14:11, Luke~18:14 and Matt.~23:12, so
the Rule attests the reception of the sentence and not demonstrably of this
pericope. Gregory the Great quotes it likewise at \work{Regula pastoralis}
III.18, where the NPNF editors refer it to Luke~18:14; his nearest citation of
this chapter is at III.21 and is of Luke~14:12--14, which the Missal does not
appoint.

\begin{witnessnote}
\wlocus{Prosper Guéranger, \work{The Liturgical Year}, vol.~XI, printed pp.~365, 367--368}
He relays Ambrose at this Gospel and says the Church has chosen him: ``she bids
them listen to St.~Ambrose, whom she has selected as her homilist for this Sunday
\dots\ he may become like the man mentioned in to-day's Gospel, who had the
dropsy; and dropsy, says our saintly preacher of Milan, is a morbid exuberance of
humours, which stupefy the soul, and induce a total extinction of spiritual
ardour. And yet, even if he were to have such a fall as that, let him not forget
that the heavenly physician is ever ready to cure him.'' \textbf{That the Breviary
appoints Ambrose here is his statement and rests on him alone.} At p.~365 he
makes the parable's wedding ``that of heaven, of which there is a prelude given
here below, by the union effected in the sacred banquet of holy Communion,''
reading it through Matt.~22:2 rather than through Luke; and at p.~368 he joins
the Gospel's closing axiom directly to grace through Ecclus.~3:20--21: ``The
greater thou art, the more humble thyself in all things, and thou shalt find
grace before God; for great is the power of God alone, and He is honoured by the
humble.'' \textbf{No other witness on this formulary as a whole
joins two of its elements outright.}
\end{witnessnote}

Two further canonical relations belong to the appointed words. \latin{Amice}
occurs in exactly five verses of the Clementine, and Luke~14:10 and
Matt.~22:12 are the only two where a host at a wedding feast addresses a guest
so --- in the second the guest is expelled. And the logion stands in a family
with Matt.~23:12, Luke~1:52, Jas.~4:10 and 1 Pet.~5:6, and sapientially with
Prov.~25:6--7 and Ecclus.~3:20--21, of which Prov.~25:7's \latin{Ascende huc} is
the nearest Old Testament wording to \latin{Amíce, ascénde supérius}, in the same
situation. \textbf{The appointed Latin is one clause shorter than the
Greek-based English at 14:3}: the missal and the Clementine read \latin{Si licet
sábbato curáre?} where the Revised Version prints ``or not.''
\latin{Legisperítus}, the word the pericope gives the Lord's interlocutors,
stands seven times in Luke and once elsewhere in the whole Clementine, at
Tit.~3:13, and Luke~14:3 is its last occurrence in that Gospel; the other
synoptic sabbath controversies name the objectors otherwise.

\subsection*{7. Offertory \cue{Off.} --- \latin{Dómine, in auxílium meum réspice}}

The antiphon takes the second half of v.~14, then the first half of v.~15, then
repeats the second half of v.~14. \textbf{The repetition is printed in full and
is not a cue}, and the psalm has no such refrain. It is the best attested of the
five chants and the one where the witnesses divide most sharply on what the psalm
is doing.

\begin{witnessnote}
\wlocus{Augustine, \work{Enarratio in Ps.~39} \S\S23--25} At v.~14 the heart
that has forsaken me is the heart that cannot comprehend itself --- Peter at
Matt.~26:35 --- and the cry is that of \latin{membra sub ferramentis medici
clamantia, sed sperantia}. \S24 states the speaker: \latin{Christus in passione
loquitur.} At v.~15 the confounding of those who sought the soul is fulfilled in
the Resurrection, \latin{illi gavisi sunt cum posuit, confusi sunt cum recepit}
(John~10:18). At the clause the Offertory does not print, \latin{convertantur
retrorsum}, he reads benevolently: that the proud who tried to walk in front of
their Lord be turned into followers behind him. \textbf{His lemma is a third
form of the petition}, \latin{Placeat tibi \dots\ Domine, in adjuvandum mihi
respice}, neither the missal's nor the Clementine's.
\end{witnessnote}

\begin{witnessnote}
\wlocus{Cassiodorus, \work{Expositio psalmorum}, in Ps.~XXXIX \S\S21--24}
\textbf{His lemma is the missal's clause word for word}: \latin{Complaceat tibi,
Domine, ut eripias me: Domine, in auxilium meum respice}, with no manuscript
variant recorded against it in Migne's apparatus. His gloss on it is the thesis
Aquinas states seven centuries later: \latin{ut intelligamus respectum ipsius
nostrum esse praesidium; sicut est illud Evangelii: Respexit Petrum, et flevit
amare. Aliter enim liberari non possumus, nisi nos Divinitas propitiata
respiciat.} On the imprecation he reads conversion --- \latin{Confundantur,
dixit, mirabilium operatione turbentur. Revereantur autem, resurrectionis gloria
corrigantur \dots\ sicut persecuti sunt, ita et praedestinati conversionis munere
liberentur} --- and distinguishes two ways of seeking a soul, \latin{sive ad
honorem, sive ad mortem}, the added \latin{ut auferant eam} fixing which is
meant. In his psalter the clause at Ps.~70:12 reads \latin{in adiutorium meum
respice} and Ps.~39:14 \latin{in auxilium meum respice}, which is the reverse of
the Clementine's distribution.
\end{witnessnote}

\begin{witnessnote}
\wlocus{Thomas Aquinas, \work{Postilla super Psalmos} in ps.~39 nn.~6--7} The
one appointed psalm of this Mass his commentary reaches. \latin{Respectus Dei
est auxilium nostrum} answers the petition directly. He allows two confusions,
of penitence (Rom.~6:21) and of punishment; he takes \latin{convertantur
retrorsum} \latin{in bonum, idest sequantur Christum retro eum}, which is
Augustine's and Cassiodorus's reading; and at \latin{cor meum dereliquit me} he
reads the distraction of fervour by venial sins, citing 2 Kings (2 Samuel) 7:27
--- \textbf{the same verse Augustine reaches for at this Mass's Introit}, for the
same fugitive heart.
\end{witnessnote}

\begin{witnessnote}
\wlocus{Robert Bellarmine, \work{Explanatio in Psalmos}, printed p.~122, his own vv.~13--14}
\textbf{He breaks with the other Latins twice.} He reads the first verse as the
Passion --- ``you seem as if you had for some time abandoned me \dots\ but now
`look down to help me,' that you may at once replenish me in the joy of a
glorious resurrection'' --- and calls the imprecation a prophecy in disguise,
``That he now prophesies in the form of an imprecation, a thing usual with the
prophets \dots\ But immediately after, when they heard of his resurrection \dots\
`they were confounded and ashamed.''' And where Augustine, Cassiodorus and
Aquinas all read the next verse's turning backward as a prayer for conversion, he
does not: it is more of the same confusion. His remark about ``the Jews'' at this
place is sixteenth-century polemical register and is his.
\end{witnessnote}

\begin{witnessnote}
\wlocus{Theodoret of Cyrus, \work{Interpretatio in Psalmos}, PG~80 cols.\ 1159--1160}
He supplies a speaker none of the Latins gives. At v.~13, immediately before the
appointed verses and governing his reading of them: \latin{Sic igitur in hoc loco
Dei Ecclesia, impiorum fluctibus concussa, non effertur, tanquam pugnare volens,
sed quae accidunt, peccatis et delictis tribuit \dots\ nec Dei Ecclesia ex
hominibus perfectis tota constat, sed habet etiam ignavos \dots\ Et quoniam unum
est corpus, tanquam ex una persona et haec et illa proferuntur.} \textbf{That is
a different solution from Augustine's \latin{totus Christus} to the same
problem.} On the appointed verses his gloss is one sentence and it is about
intent, \latin{Non enim vulnerare volunt, sed sempiternae morti tradere}; at the
unprinted next verse he reads rout and not conversion, \latin{Verte ipsos in
fugam.}
\end{witnessnote}

\begin{witnessnote}
\wlocus{Jerome, \work{Epistula} 106 \S23, on this very verse} His correspondents
found in their Greek the verb the Latin renders \latin{festina}; he answers that
the Septuagint has the verb the Latin renders \latin{respice}, and lets the
Latin stand.
\textbf{So the verb the Offertory sings twice was contested in the fourth century
and Jerome ruled for it.} What he settles is the verb and not the noun: his own
form of the clause is \latin{in adiutorium meum respice}, a fourth Latin again,
and nothing in the letter bears on \latin{auxilium} against \latin{adiutorium}.
The variant he corrects, \latin{festina}, is the reading of Ps.~69:2, the verse
this psalm doubles.
\end{witnessnote}

\textbf{Four Latin forms of one clause stand among these witnesses}: Augustine's
\latin{in adjuvandum mihi respice}; the Clementine's and Theodoret's Migne
Latin's \latin{ad adjuvandum me respice}; Cassiodorus's and the missal's
\latin{in auxilium meum respice}; and Jerome's \latin{in adiutorium meum
respice}. \textbf{No commentary here has a counterpart to the printed refrain},
and none of them is expounding a chant.

\textbf{Not one of the five reads \latin{confundántur et revereántur} as wishing
the enemies' destruction, and the five relieve it four ways}: Augustine and
Bellarmine together as a prophecy already fulfilled at the Resurrection,
Cassiodorus as a prayer for the persecutors' conversion, Theodoret as a petition
that historical enemies be routed and fail of what they want, Aquinas by
allowing the penitential and the punitive confusion together. The gentlest of
them is a Latin reading and not the tradition's. \textbf{Guéranger, above this Offertory at
printed p.~370, is the one witness who keeps the enemies as enemies}: ``The
greater the conquests made by the Church, the greater are the efforts of hell to
destroy the souls of her dear children. This fearful danger calls for her fervent
prayers; and our Offertory-anthem is one of these.''

The chant's own reading is verbatim Ps.~70:12, a verse of this Mass's own
Communion psalm, with the subject differing. \latin{Simul} stands in Ps.~39:15
and the chant drops it, and in dropping it the chant agrees with the Ps.~69:2--6
doublet at Ps.~69:3, which does not carry the word; but the Offertory keeps
\latin{ut áuferant eam}, which Ps.~69:3 has not, so the chant is not the
Ps.~69 form and the agreement is partial. The psalm is also that of \latin{Ecce
venio \dots\ ut facerem voluntatem tuam} at vv.~7--9, quoted at Heb.~10:5--7 of
Christ entering the world --- though \textbf{those verses lie outside the
appointed excerpt}, and the Latin of the quotation differs from the Latin of the
psalm at the decisive clause, \latin{corpus autem aptasti mihi} against
\latin{aures autem perfecisti mihi}.

This antiphon's wording produced no later reuse: \latin{in auxílium meum
réspice} is close enough to Ps.~69:2's \latin{Deus in adiutorium meum intende},
the versicle that opens every Hour, that any later use of those words is a use
of that verse. The calendar registry gives this Offertory one further
appointment, at the Friday after the Second Sunday of Lent; Schuster names the
same one.

\subsection*{8. Secret \cue{Sec.} --- \latin{Munda nos, quǽsumus, Dómine}}

The prayer asks two things in one movement: that the effect of the sacrifice now
offered cleanse us, and that God mercifully complete in us what makes us fit to
share it. \textbf{\latin{Mereámur} stands in a purpose clause under
\latin{pérfice miserátus in nobis}}, so the merit named is God's completed work
and not a prior claim.

\textbf{No Father reads this prayer or the Postcommunion.} The Latin corpora
that carry the reception of this Mass's other elements --- Augustine's
\work{Enarrationes} on Pss.~31--40, 61--70, 81--90, 91--100 and 101--110 and his
\work{Quaestiones euangeliorum}, Ambrose's \work{Expositio} VII--IX, Bede's
\work{In Lucam} entire, Benedict's \work{Regula} entire, Aquinas's \work{Catena
in Lucam} 14--18, \work{Super Ephesios} entire and \work{Super Psalmos} 31--40,
and the Tridentine text --- carry neither \latin{sacrificii praesentis} nor
\latin{mereamur esse participes} nor \latin{Munda nos}. \textbf{The one witness
who reads it is Guéranger}, above this Secret at printed p.~370: ``the Sacrifice,
at which we are present, and which is to be consummated, in a few moments, by the
words of Consecration, is the most direct and efficacious of all the immediate
preparations that we can make for the Communion of the Body and Blood, which that
Sacrifice produces on the altar.''

What the older books supply is the prayer's provenance. \textbf{The Secret
and the Postcommunion stand together, word for word, in the Old Gelasian}, Book
III sect.~XII no.~693 at printed p.~231 of Wilson's 1894 edition, under the
heading \latin{ITEM ALIA MISSA} and under two collects the Roman Missal does not
use here; \latin{Tua nos} is not in that book's three books at all. All three
stand together in the Gregorian Hadrianum at sect.~XXXV, printed p.~174 of
Wilson's 1915 edition, under the heading \latin{DOMINICA .XVII. POST
PENTECOSTEN}; and all three again in the Frankish Gelasian at printed p.~184 of
Gerbert's \work{Monumenta veteris liturgiae Alemannicae}, under \latin{Dominica
XX. Post Pentecosten. Gelas.}, with a proper preface \latin{VD. \dots\ Precantes,
ut Ihs Xps Filius tuus Dns noster sua nos gratia protegat}. \textbf{Gerbert marks
each oration}: \latin{Tua nos} \latin{gg.}, Gregorian as Muratori has it, alone;
the Secret and the Postcommunion \latin{g. El. gg.}, Gelasian, Eligian and
Gregorian together. \textbf{The Veronense has none of the three}; its one near
miss, \latin{tua nos semper gratia praeveniens laetificatur}, is a different
prayer.

The 1862 Pustet printed \latin{Alia Secreta} beside this one, and the Vatican
typica of 1604 carries \latin{alia Secreta} and \latin{tertia ad libitum} in the
same window. Those belong to the discipline the general decree of 23~March 1955
abolished and are no part of this formulary.

\subsection*{9. Communion \cue{Comm.} --- \latin{Dómine, memorábor iustítiæ tuæ solíus}}

The exposition that stands over this antiphon is an exposition of exactly its
three verses and of nothing else. Four commentators reach them, and they divide
at both of the antiphon's hinges --- what \latin{solíus} excludes, and what the
three ages are.

\begin{witnessnote}
\wlocus{Augustine, \work{Enarratio in Ps.~70}, sermo~II \S\S1, 2 and 4} The
three sections are Ps.~70:16, 17 and 18, the antiphon's own extent. On v.~16 he
fastens on the one word: \latin{Domine, memorabor iustitiae tuae solius. O
solius! Quid addidit, solius? \dots\ Justitia tua sola me liberat; mea sola non
sunt nisi peccata} (1 Cor.~4:7); and \latin{Gratia gratis data est: nam nisi
gratis esset, gratia non esset \dots\ nihil tuum praecessit, ut acciperes.} On
v.~17: \latin{Quid me docuisti? Quia tuae solius iustitiae memorari debeo \dots\
Debebatur poena; reddita est gratia \dots\ docuisti me nihil in me
praecessisse.} On v.~18 he notes the Greek pair behind the Latin doublet
\latin{senéctam et sénium}, then gives two senses: individually, \latin{usque ad
ultimum meum, nisi mecum fueris, non erit aliquid meriti mei: gratia tua semper
perseveret mecum}; and ecclesially, \latin{vox est enim Ecclesiae}, the Church
whose youth was the age of the apostles and who endures to the end against those
who predicted that Christianity would last a season. \S3 of the same discourse
cites \textbf{Ps.~85:11}, the psalm of this Mass's own Introit, for the point
that grace must lead as well as start.
\end{witnessnote}

\begin{witnessnote}
\wlocus{Cassiodorus, \work{Expositio psalmorum}, in Ps.~LXX \S\S4, 21--23, 32--33}
On the psalm as a whole he says what Augustine says and nearly in Augustine's
words: the speaker preaches Christ's charity \latin{quae nullis meritis
praecedentibus gratis semper impenditur}, and \latin{Cum totus hic psalmus
gratiam Domini, quae gratis datur, summa intentione commendet}. \textbf{On the
antiphon's first clause he disagrees with him.} \latin{Solius} does not exclude
the singer's own righteousness but refers the remembering to the Judgment ---
\latin{illo scilicet tempore cum agnos sequestrat ab haedis \dots\ Tunc enim vere
memor erit iustitiae solius Domini, quoniam eam et mirabilem, et singularem esse
cognoscit.} On \latin{a iuventute mea} youth is the age at which a man first
comes to grace; on \latin{usque in senectam et senium} he makes Augustine's
philological observation independently and then periodises the Church's ages
differently, her youth the Crucifixion and the martyrs, her \latin{senecta} the
present age near the end, her \latin{senium} the last declining time
\latin{quando et saevus ille tyrannus adveniet}. \textbf{His Judgment gloss here
is the same sheep-and-goats formula he uses at Ps.~101:17 for the Gradual's
\latin{videbitur in maiestate sua}}, so in this one commentator the Gradual's
second half and the Communion's first clause name one event.
\end{witnessnote}

\begin{witnessnote}
\wlocus{Robert Bellarmine, \work{Explanatio in Psalmos}, printed pp.~212--213}
The literal-historical voice of the four. On \latin{solíus} the exclusion is of
every other support --- ``I will lose sight completely of human counsel, of my
own strength, or of my friends \dots\ `thy justice alone,' by virtue of which you
keep your promises.'' On \latin{a iuventute mea} he declines allegory: ``it was
in consequence, that I, an unarmed youth, fought with a bear and a lion, and
conquered both them and the giant Goliath.'' And he makes the petition run to the
completion of the psalter, ``until I shall have finished the book of Psalms,
through which I will shew forth thy arm to all posterity'' --- a striking reading
for a Communion antiphon, and the clause carrying it stands outside the appointed
text.
\end{witnessnote}

\begin{witnessnote}
\wlocus{Theodoret of Cyrus, \work{Interpretatio in Psalmos}, PG~80 cols.\ 1423--1426}
\textbf{His periodisation is unlike anything the Latins give: youth is the age of
Moses and the giving of the Law, and old age is the Law's end.} \latin{Atque hinc
manifestum est, eum magni Mosis tempus juventutem vocare. Per illum enim lex data
fuit}; \latin{Prophetice David futura praedicit, et senectutem vocat legis finem.
Hoc autem post Domini Christi adventum contigit} --- with Heb.~8:13 in support,
\latin{quod autem antiquatur et senescit, prope interitum est}. On
\latin{memorabor iustitiae tuae} he anchors the clause in the preceding verse's
justice, the just verdict by which God judged between the speaker and his
enemies. And he cross-refers \latin{quoniam non cognovi litteraturam} to
Ps.~39:5--6, a verse of this Mass's own Offertory psalm, as Migne's footnote
records. In his Greek the word for \emph{alone} stands in the
genitive and may attach to \emph{thy} rather than to \emph{justice}, where the
Latin \latin{iustitiae tuae solius} attaches it to the justice; Augustine's whole
reading of the clause rests on the Latin, and the facsimile does not settle the
Greek construal.
\end{witnessnote}

\begin{witnessnote}
\wlocus{Jerome, \work{Epistula} 106 \S43, on two clauses of these verses}
Against Greek copies reading \latin{Deus meus} he holds that \latin{meus} is
superfluous in \latin{Deus, docuisti me a iuventute mea}; and against copies
reading \latin{mirabilia tua} at \latin{donec annuntiem brachium tuum} he holds
that \latin{mirabilia tua} has been carried up from the preceding verse and that
\latin{brachium} is right. \textbf{He is adjudicating the exact wording this
antiphon sings.}
\end{witnessnote}

\textbf{So the antiphon's first clause carries four readings and its span
carries four more.} \latin{Solíus} excludes, for Augustine, any righteousness of
the singer's own; for Bellarmine, human counsel and his own strength; for
Cassiodorus it refers the remembering to the Judgment; for Theodoret it is God's
just dealing with the speaker. \textbf{The grace-alone sense is Augustine's, and
Bellarmine's in a different key, and it is not the undivided tradition's.} The
three ages are Augustine's beginning of faith and, ecclesially, the age of the
apostles; Cassiodorus's coming to grace, with the last age Antichrist's;
Bellarmine's David's own boyhood; Theodoret's the age of Moses.
\textbf{All four read the span as a span of God's teaching and none as a lifetime
in the ordinary sense.} And Guéranger, above this antiphon at printed
pp.~370--371, reads it from outside the exegetical tradition altogether: the
Church, filled with the Wisdom of the Father, ``promises God, as her
thank-offering, that she will keep His justice, which is His law, and that she
will labour to make His divine teaching produce its fruits.''

\textbf{Augustine's most quotable formula on this doctrine, \latin{sua dona
coronabit, non merita tua}, stands in sermo~II \S5, on v.~19} --- one verse past
where the antiphon stops --- so it is not a comment on the sung text. And the
NPNF English runs a single section series across both discourses where the Migne
Latin restarts at 1 for Sermo~II.

The antiphon is the only appointed scriptural element of this Mass with no
verbal departure from the Clementine, and it restates the psalm's own earlier
petitions at vv.~5 and 9. Its rare genitive \latin{solíus} stands in four verses
of the whole Clementine, and one of them is \textbf{Ecclus.~3:21},
\latin{magna potentia Dei solius, et ab humilibus honoratur}, inside the
sapiential humility passage behind the appointed Gospel's conclusion and the
verse Guéranger reaches for there. \textbf{Only \latin{solíus} is inside
appointed text}: \latin{potentias Domini} is the unsung first half of v.~16.
\latin{Sénium} occurs nowhere else in the Clementine.

The calendar registry gives the Communion one further appointment, the Thursday
after the Fourth Sunday of Lent, which Schuster names; the 1962 book's own text
layer adds a second, a votive Mass \latin{pro quacumque necessitate} with
\latin{T. P. Alleluia} added, the registry holding no votive Masses at all.
\textbf{That
Lenten feria is also where the Fifteenth Sunday's Gospel is read}, Luke~7:11--16,
so one Lenten Mass holds together what the two September Sundays hold apart, and
Lasance's 1945 hand missal prints the same pairing.

The Douay's wording of these verses has no later use: ``grey hairs'' stands in
English and American literature as an ageing character described and nothing
else, and ``I will be mindful of thy justice'' returns no newspaper page in
\work{Chronicling America}. The Authorized Version, searched for both, prints
neither at these verses.

\textbf{No New Testament book quotes any verse this Mass appoints.} Heb.~1:10--12
quotes Ps.~101 and Heb.~10:5--7 quotes Ps.~39, both outside the excerpts the
chants sing, and Ps.~85:1, 3 and 5, Ps.~97:1, Ps.~39:14--15 and Ps.~70:16--18 are
quoted in no book from Matthew to the Apocalypse. \textbf{Every christological
reading of these five chants therefore rests on the psalms' wider context, on
typology, or on the commentators above, and on no apostolic citation of the
words actually sung.}

\subsection*{10. Postcommunion \cue{Postcomm.} --- \latin{Purífica, quǽsumus, Dómine, mentes nostras}}

The last oration asks for cleansing and renewal by the sacraments just received,
and then for help to the body, \latin{præsens páriter et futúrum}. \textbf{Two
words tie it back into the formulary}: \latin{auxílium}, which stands in this
Mass at the Offertory and here and nowhere else, and \latin{rénova}, which asks
as a petition the renewal the Epistle asks for \latin{in interiórem hóminem} and
puts inside it the newness the Alleluia sings as \latin{cánticum novum}.

\textbf{How the Secret's \latin{sacrifícii præséntis efféctu} stands to this
prayer's \latin{rénova cæléstibus sacraméntis} --- offered sacrifice against
received sacrament --- neither text says.} Guéranger gives a direction of
dependence and only he: the Sacrifice is the immediate preparation for the
Communion it produces, so the two orations are two moments of one movement. And
above this prayer, at printed p.~371, he construes \latin{præsens páriter et
futúrum} temporally: ``let us pray, with the Church, that we may be renewed by
the purity, which these heavenly mysteries bring to us, who are well prepared for
the gift: the effect of such a gift tells upon our bodies, both in this and in
the next life.'' The Latin leaves one thing he does not consider --- whether the
Secret's \latin{eius} refers to the sacrifice or to its effect.

Its provenance is the Secret's, the two standing together in the Old Gelasian
under collects this Missal does not use here.

\textbf{The 1862 Pustet corroborates this oration word for word and whole}, the
conjunction \latin{et} before \latin{rénova} included, at leaf n426, printed
p.~341: \latin{Purífica, quæsumus Dómine, mentes nostras benígnus, et rénova
cœléstibus sacraméntis: ut consequénter et córporum præsens páriter, et futúrum
capiámus auxílium. Per Dóminum.} That conjunction stands at the end of a printed
line, where this printing's optical layer loses material, and it is printed on
the page image. Two older books carry it at the same prayer: the tracked full
texts of the \work{Missale Romanum} of Venice, 1570 and of the Vatican typical
edition of 1604 both have \latin{et} before \latin{rénova}, read there at
optical state and for that word alone. The differences from the typical
edition --- \latin{cœ} for \latin{cæ}, the comma after \latin{páriter} and none
after \latin{quæsumus} --- change no word a reader would say aloud.

The 1862 book also printed \latin{Alia Postcommunio. Mundet et muniat.} and
\latin{Tertia ad libitum.} at this Sunday; both belong to the pre-1955
discipline and are no part of this formulary.

\clearpage
